% -------------------------------------------------------
\section{Conclusion}
% -------------------------------------------------------

This study examined how discussions of AI, ML, and NLP research changed over 30 years by analyzing 24,677 arXiv abstracts using five text analysis methods. These methods together provided a clear and consistent perspective, providing the basis for the findings discussed below.

The clearest finding is that 2014 to 2016 represented a real turning point, shown independently by word analysis, topic grouping, classification, tracking term spread, and data simplification methods. Notably, a strong link of 0.863 between TF-IDF and Sentence-BERT distance measures, which work at very different language levels, is the best proof that this change is a real part of the field’s structure, not just a result of the methods used. This change is better described as an increase rather than a replacement: theory-focused topics became less common but remained present, and most tracked terms shifted in meaning rather than disappearing. Building on these results,

The analysis of how terms spread gave the most new insight. Specifically, combining changes in TF-IDF, shifts in Word2Vec word groups, and changes in SPECTER2 discourse showed different ways terms were adopted that no single method could see. The main pattern, quiet shifts in meaning without matching increases in use, explained 13 of 36 tracked terms. The encoder example showed this clearly, with a 63-fold increase in mentions and a complete change in its role, visible only through SPECTER2. The pre-training example showed that copying templates is a different kind of change, which a simple frequency count would have missed. Looking at how limitations were described added another view: ethics-related words grew from almost none before 2015 to 13\% of abstracts by 2020–2024, and failure-related words became the second most common way to describe limitations, changes not visible in simple counts of limitation words. All methods performed as intended. The limitation of the framing classifier was that it measured broad subfield vocabulary shifts rather than framing-specific change, and TF-IDF PCA recovered subfield composition variance rather than temporal drift. Both failures were diagnostically informative rather than simply negative, demonstrating the value of cross-method checking. Turning to the remaining challenges,

Three difficulties remain. The SPECTER2 coverage limits exclude the newest terms, such as transformer, prompt, and downstream task, by design. The 6.9\% Procrustes anchor basis can make some findings about term changes unclear because of alignment issues. Looking only at abstracts may hide the full variety of how limitations are described in full papers. The bigger lesson is that no single method fully captures how scientific language changes. The strength of this study lies in the concurrence across methods, and their disagreements were also helpful.